%=============================================================
% Madhumitha Murthy — ATS-Friendly CV
% Compiler : XeLaTeX or pdfLaTeX (both work on Overleaf)
% Engine   : Use "XeLaTeX" in Overleaf for best font rendering
%=============================================================

\documentclass[10pt, a4paper]{article}

% ── Packages ─────────────────────────────────────────────────
\usepackage[T1]{fontenc}
\usepackage[utf8]{inputenc}
\usepackage{lmodern}
\usepackage[margin=0.75in, top=0.65in, bottom=0.65in]{geometry}
\usepackage{hyperref}
\usepackage{enumitem}
\usepackage{titlesec}
\usepackage{xcolor}
\usepackage{array}
\usepackage{tabularx}
\usepackage{multicol}
\usepackage{parskip}

% ── Colours ──────────────────────────────────────────────────
\definecolor{headingblue}{RGB}{30, 60, 110}
\definecolor{rulegray}{RGB}{180, 180, 180}

% ── Hyperlinks ───────────────────────────────────────────────
\hypersetup{
    colorlinks=true,
    urlcolor=headingblue,
    linkcolor=headingblue,
    pdfauthor={Madhumitha Murthy},
    pdftitle={Madhumitha Murthy - CV},
    pdfsubject={Machine Learning Engineer CV}
}

% ── Section formatting ────────────────────────────────────────
\titleformat{\section}
    {\large\bfseries\color{headingblue}}
    {}{0em}{}
    [\color{rulegray}\titlerule]

\titlespacing{\section}{0pt}{8pt}{4pt}

% ── List settings ─────────────────────────────────────────────
\setlist[itemize]{
    leftmargin=1.2em,
    itemsep=1pt,
    parsep=0pt,
    topsep=2pt,
    label=\textbullet
}

% ── Custom commands ───────────────────────────────────────────

% Job entry: {Title}{Company}{Dates}
\newcommand{\jobentry}[3]{%
    \noindent
    \begin{tabularx}{\textwidth}{@{}X r@{}}
        \textbf{#1} $\cdot$ #2 & \textit{#3}
    \end{tabularx}%
    \vspace{1pt}%
}

% Education entry: {Degree}{Institution}{Dates}
\newcommand{\eduentry}[3]{%
    \noindent
    \begin{tabularx}{\textwidth}{@{}X r@{}}
        \textbf{#1} & \textit{#3} \\
        #2 &
    \end{tabularx}%
    \vspace{1pt}%
}

% Project entry: {Title}{Stack summary}{Date}
\newcommand{\projentry}[3]{%
    \noindent
    \begin{tabularx}{\textwidth}{@{}X r@{}}
        \textbf{#1} \textit{\small | #2} & \textit{#3}
    \end{tabularx}%
    \vspace{1pt}%
}

% ── Spacing helpers ───────────────────────────────────────────
\newcommand{\entryspace}{\vspace{6pt}}
\newcommand{\sectionspace}{\vspace{2pt}}

%=============================================================
\begin{document}
\pagestyle{empty}

%-------------------------------------------------------------
% HEADER
%-------------------------------------------------------------
\begin{center}
    {\LARGE\bfseries\color{headingblue} Madhumitha Murthy}\\[4pt]
    {\small
        +65 8319 1000 \quad|\quad
        \href{mailto:madhumit007@e.ntu.edu.sg}{madhumit007@e.ntu.edu.sg} \quad|\quad
        \href{https://linkedin.com/in/madhumitha-murthy-4801b7223}{linkedin.com/in/madhumitha-murthy} \quad|\quad
        \href{https://github.com/madhumitha-murthy}{github.com/madhumitha-murthy} \quad|\quad
        \href{https://scholar.google.com/citations?user=4jVqyt0AAAAJ}{Google Scholar} \quad|\quad
        Singapore
    }
\end{center}

\vspace{4pt}

%-------------------------------------------------------------
% SUMMARY
%-------------------------------------------------------------
\section{Summary}
\sectionspace

Machine Learning Engineer and researcher with peer-reviewed publications at IEEE ISCAS and npj Wireless Technology, specialising in ISAC systems, 5G/6G signal processing, and multilingual NLP. At A*STAR I2R, built a high-fidelity 3D ray-tracing simulation platform using NVIDIA Sionna for vehicle sensing and communication in urban environments, proposing novel algorithms for clutter cancellation, trajectory prediction, and sensing-aided beamforming. Experienced in end-to-end ML pipelines, LLM fine-tuning, RAG systems, time-series anomaly detection, and production deployment with FastAPI, Docker, and MLflow. Graduating April 2026 from NTU (MSc, EEE). Open to Machine Learning Engineer, NLP Engineer, and Data Scientist roles in Singapore.

\entryspace

%-------------------------------------------------------------
% EDUCATION
%-------------------------------------------------------------
\section{Education}
\sectionspace

\eduentry{MSc in Signal Processing and Machine Learning}{Nanyang Technological University (NTU), School of EEE, Singapore}{Aug 2024 -- Apr 2026 (Expected)}
\begin{itemize}
    \item Dissertation with I2R, A*STAR on ISAC using 5G for V2V/Infrastructure Communication and Localisation (Jan -- Apr 2026: full-time dissertation)
    \item Coursework: Genetic Algorithms \& ML, NLP, AI \& Data Mining, Pattern Recognition \& Deep Learning, Probability \& Random Processes, Real-Time DSP, Machine Vision
\end{itemize}

\entryspace

\eduentry{Bachelor of Engineering in Computer Science and Engineering}{Meenakshi Sundararajan Engineering College, Anna University, Chennai, India}{Sept 2020 -- July 2024}
\begin{itemize}
    \item First-class honours. Coursework: Data Structures \& Algorithms, Software Engineering, Database Systems, Computer Networks.
\end{itemize}

\entryspace

%-------------------------------------------------------------
% EXPERIENCE
%-------------------------------------------------------------
\section{Experience}
\sectionspace

\jobentry{Research Attachment}{Institute for Infocomm Research (I2R), A*STAR, Singapore}{Oct 2024 -- Present}
\begin{itemize}
    \item Built a high-fidelity 3D ISAC simulation platform using NVIDIA Sionna ray tracing, modelling multipath propagation, Doppler effects, and dynamic occlusions for vehicle sensing and communication in a realistic urban environment (Fusionopolis, Singapore).
    \item Constructed a 3D digital twin by integrating OpenStreetMap road geometry and Blender-modelled building infrastructure, enabling physics-consistent signal interaction modelling across monostatic and bistatic radar configurations.
    \item Designed and implemented an OFDM radar algorithm with Static Clutter Cancellation (SCC) at 59 GHz / 100 MHz bandwidth, suppressing static reflections to isolate moving vehicle targets with distinct Doppler signatures.
    \item Developed a Doppler-based vehicle trajectory prediction algorithm forecasting future vehicle positions from past sensing estimates, achieving close trajectory tracking at update intervals as short as 0.333s --- enabling proactive beam alignment in high-mobility scenarios.
    \item Implemented sensing-aided beamforming leveraging predicted vehicle positions to update beamforming weights in real time, reducing computational complexity versus true DOA-based benchmarks.
    \item Co-authored two peer-reviewed publications: accepted at IEEE ISCAS 2025 and npj Wireless Technology.
\end{itemize}

\entryspace

\jobentry{Research Assistant}{NTU, College of Computing \& Data Science, Singapore}{Sept 2025 -- Dec 2025}
\begin{itemize}
    \item Conducted literature review on rubric-based peer evaluation methodologies and developed specifications for a bias-resilient peer evaluation system for NTU.
\end{itemize}

\entryspace

\jobentry{Postgraduate Teaching Assistant}{NTU, School of EEE, Singapore}{Aug 2024 -- Jan 2025}
\begin{itemize}
    \item Developed lab manual and reference implementation for object detection and instance segmentation using Mask R-CNN; graded assignments for IE2108: Data Structures \& Algorithms in Python (two cohort groups).
\end{itemize}

\entryspace

\jobentry{Android Project Trainee}{Zoho Corporation, Chennai, India}{Feb 2024 -- Mar 2024}
\begin{itemize}
    \item Built Kotlin-based banking application with account management and transaction features using Coroutines and asynchronous workflows.
    \item Won 1st place at Cliq-Trix 2024 (internal bot-building competition); awarded Rs.\,1,00,000 prize and received internship offer.
\end{itemize}

\entryspace

%-------------------------------------------------------------
% PROJECTS
%-------------------------------------------------------------
\section{Projects}
\sectionspace

%% ── Project 1: RAG ──────────────────────────────────────────
\projentry{RAG-Powered Document Q\&A API}{LangChain · FAISS · HuggingFace · FastAPI · Docker · MLflow}{2025}
\begin{itemize}
    \item Built an end-to-end Retrieval-Augmented Generation (RAG) pipeline enabling users to upload any PDF and receive grounded, context-specific answers via a REST API.
    \item Chunked and embedded documents using HuggingFace sentence-transformers; stored and retrieved vectors using FAISS for sub-second semantic search.
    \item Wrapped the pipeline in a FastAPI application with \texttt{/upload} and \texttt{/query} endpoints; containerised with Docker and deployed on HuggingFace Spaces.
    \item Tracked 3 retrieval configurations (varying chunk size and top-$k$) using MLflow, comparing latency and answer quality metrics across runs.
    \item \textit{GitHub:} \href{https://github.com/madhumitha-murthy/rag-document-qa}{github.com/madhumitha-murthy/rag-document-qa}
\end{itemize}

\entryspace

%% ── Project 2: Fine-tuning ──────────────────────────────────
\projentry{Domain-Specific LLM Fine-Tuning with LoRA}{HuggingFace PEFT · LoRA · DistilBERT · FastAPI · Docker · MLflow}{2025}
\begin{itemize}
    \item Fine-tuned DistilBERT on the FinancialPhraseBank dataset for financial sentiment classification using parameter-efficient fine-tuning (LoRA via HuggingFace PEFT), reducing trainable parameters by over 90\% versus full fine-tuning.
    \item Tracked 3 training runs with MLflow (varying learning rate and LoRA rank), logging training loss, validation accuracy, and inference latency; selected best checkpoint by F1 score.
    \item Served the model via a FastAPI inference endpoint containerised with Docker.
    \item \textit{GitHub:} \href{https://github.com/madhumitha-murthy/llm-finetuning-lora}{github.com/madhumitha-murthy/llm-finetuning-lora}
\end{itemize}

\entryspace

%% ── Project 3: Time-series ──────────────────────────────────
\projentry{Time-Series Anomaly Detection Pipeline}{PyTorch · LSTM Autoencoder · scikit-learn · FastAPI · Docker · MLflow}{2025}
\begin{itemize}
    \item Built an end-to-end anomaly detection pipeline on the NASA SMAP satellite sensor dataset, comparing a statistical baseline (Isolation Forest) against a deep learning model (LSTM Autoencoder).
    \item LSTM Autoencoder learned normal sensor behaviour from multivariate time-series windows; anomalies flagged when reconstruction error exceeded a threshold tuned on the validation set.
    \item Logged both models to MLflow with precision, recall, F1, and inference latency; compared performance across retrieval configurations.
    \item Deployed a FastAPI endpoint accepting a time-series window and returning anomaly score and binary flag; containerised with Docker.
    \item \textit{GitHub:} \href{https://github.com/madhumitha-murthy/timeseries-anomaly-detection}{github.com/madhumitha-murthy/timeseries-anomaly-detection}
\end{itemize}

\entryspace

%% ── Project 4: Fake News ────────────────────────────────────
\projentry{Fake News Detection in Dravidian Languages}{XLM-RoBERTa · mBERT · ALBERT · HuggingFace}{Mar 2024}
\begin{itemize}
    \item Fine-tuned XLM-RoBERTa, mBERT, and ALBERT on multilingual social media datasets for misinformation detection; achieved Macro-F1 of \textbf{0.84} and \textbf{3rd place} on the DravidianLangTech 2024 shared task leaderboard.
    \item Built end-to-end NLP pipeline covering multilingual tokenisation, transformer fine-tuning, and evaluation using Python and HuggingFace Transformers.
    \item Published at \href{https://aclanthology.org/2024.dravidianlangtech-1.0}{DravidianLangTech Workshop, EACL 2024}.
\end{itemize}

\entryspace

%% ── Project 5: Sarcasm ──────────────────────────────────────
\projentry{Sarcasm Identification in Dravidian Languages}{mBERT · ALBERT · HuggingFace}{Aug 2023}
\begin{itemize}
    \item Developed transformer-based models for sarcasm detection in code-mixed Dravidian social media text; achieved Macro-F1 of \textbf{0.70} and \textbf{0.68}, ranking \textbf{8th and 5th} on two official FIRE 2023 leaderboards.
    \item Published in \href{https://ceur-ws.org}{CEUR Workshop Proceedings, FIRE 2023}.
\end{itemize}

\entryspace

%% ── Project 6: Depression ───────────────────────────────────
\projentry{Depression Detection in Social Media Text}{RoBERTa · ALBERT · HuggingFace}{Jun 2023}
\begin{itemize}
    \item Fine-tuned ALBERT and RoBERTa to classify depression severity levels from social media text; implemented complete NLP pipeline from preprocessing to evaluation.
    \item Published at \href{https://aclanthology.org/2023.ltedi-1.0}{LT-EDI Workshop, ACL 2023}.
\end{itemize}

\entryspace

%% ── Project 7: Object Detection ─────────────────────────────
\projentry{Object Detection Lab --- Mask R-CNN}{PyTorch · Mask R-CNN · Computer Vision}{Jan 2025}
\begin{itemize}
    \item Developed a complete lab manual and reference implementation for instance segmentation using Mask R-CNN, adopted in NTU's postgraduate Machine Vision curriculum.
    \item Covered end-to-end pipeline: dataset preparation, anchor configuration, region proposal network, mask head training, and evaluation using COCO metrics.
\end{itemize}

\entryspace

%-------------------------------------------------------------
% PUBLICATIONS
%-------------------------------------------------------------
\section{Publications}
\sectionspace

\begin{itemize}
    \item \textbf{IEEE ISCAS 2025} (Accepted, Presented) --- ISAC for Intelligent Transportation: Ray-Tracing, Clutter Cancellation, and Sensing-Aided Beamforming.
    \item \textbf{npj Wireless Technology} (Accepted) --- Integrated Communication and Sensing: Algorithms \& 3D Simulation Insights.
    \item \textbf{EACL 2024}, DravidianLangTech Workshop --- \href{https://aclanthology.org/2024.dravidianlangtech-1.0}{TechWhiz@DravidianLangTech 2024: Fake News Detection Using Deep Learning Models}.
    \item \textbf{CEUR Workshop Proceedings, FIRE 2023} --- \href{https://ceur-ws.org}{Sarcasm Detection in Dravidian Languages using Transformer Models}.
    \item \textbf{ACL 2023}, LT-EDI Workshop --- \href{https://aclanthology.org/2023.ltedi-1.0}{TechWhiz@LT-EDI-2023: Transformer Models to Detect Levels of Depression from Social Media Text}.
\end{itemize}

\entryspace

%-------------------------------------------------------------
% SKILLS
%-------------------------------------------------------------
\section{Skills}
\sectionspace

\begin{tabularx}{\textwidth}{@{}>{\bfseries}l X@{}}
    Programming        & Python, MATLAB, SQL, Kotlin \\[2pt]
    ML / Deep Learning & PyTorch, TensorFlow, HuggingFace Transformers, scikit-learn, LLMs, LSTM \\[2pt]
    NLP                & Transformer fine-tuning, LoRA / PEFT, multilingual NLP, RAG, FAISS, XLM-RoBERTa, mBERT \\[2pt]
    Time-Series        & Anomaly detection, LSTM Autoencoder, Isolation Forest, multivariate forecasting \\[2pt]
    Signal Processing  & OFDM Radar, 5G/6G NR, ISAC, Beamforming, Delay-Doppler Processing, Channel Modelling \\[2pt]
    Simulation         & NVIDIA Sionna ray tracing, Blender, OpenStreetMap, digital twin \\[2pt]
    MLOps / Deploy     & FastAPI, Docker, MLflow, Git, HuggingFace Spaces \\[2pt]
    Scientific         & NumPy, SciPy \\
\end{tabularx}

\entryspace

%-------------------------------------------------------------
% ACHIEVEMENTS
%-------------------------------------------------------------
\section{Achievements}
\sectionspace

\begin{itemize}
    \item \textbf{1st Place}, CliqTrix 2024 (Zoho Corporation) --- Bot-Building Competition. Awarded Rs.\,1,00,000 prize and internship offer.
    \item \textbf{3rd Place}, DravidianLangTech 2024 --- Fake News Detection Shared Task (Macro-F1: 0.84).
    \item \textbf{Top-10}, FIRE 2023 --- Sarcasm Detection in Dravidian Languages (8th and 5th place across two tasks).
\end{itemize}

\end{document}
